\subsection{Complete period collapse can hide arbitrary parameter degrees}

\begin{proposition}[P-F21: no quantitative transfer from parameters to positive forms in general]
\label{prop:toricperiodcollapse}
For every integer $D\ge2$, there is a positively graded complex domain
$R_D$ of Krull dimension three with the following properties:
\begin{enumerate}
\item $R_D$ is normal, Cohen--Macaulay, and has rational singularities.
\item Its global Hilbert function and Hilbert series are
\[
 f_D(s)=\frac{D-1}{2}s^2+\frac{D+1}{2}s+1\quad(s\ge0),
 \qquad
 \operatorname{Hilb}_{R_D}(t)=\frac{1+(D-2)t}{(1-t)^3}.
 \tag{P-F21a}
\]
Thus it has a positive numerator and three denominator weights equal to
one, and its positivity and saturation indices are both zero.
\item Every homogeneous system of parameters has at least one degree
divisible by $D$. There is a system of parameter degrees $(1,1,D)$.
Consequently the maximal degree in the lexicographically minimal
nondecreasing parameter-degree tuple is exactly $D$.
\item It has a homogeneous quotient $\mathbb C[z]$ with $\deg z=D$.
In particular its universal Hilbert denominator contains $\Phi_D$,
although its own reduced Hilbert denominator has no nontrivial root
of unity as a zero.
\end{enumerate}
This statement concerns general graded domains; no identification of
$R_D$ with an ordinary irreducible-inner plethysm algebra is asserted.
\end{proposition}

\begin{proof}
Consider the rational triangle
\[
 P_D=\operatorname{conv}\left\{(0,0),(D,0),
           \left(0,\frac{D-1}{D}\right)\right\}
\]
and its integral cone semigroup
\[
 S_D=\{(x,y,s)\in\mathbb Z^3:s\ge0,\ (x,y)\in sP_D\},
 \qquad R_D=\mathbb C[S_D].
\]
The monomial corresponding to $v=(x,y,s)$ is denoted $X^v$ and
has degree $s$. The degree-zero part is $\mathbb C$, and every graded
part is finite-dimensional. The cone has linearly independent
primitive rays
\[
 v_0=(0,0,1),\qquad v_1=(D,0,1),\qquad v_2=(0,D-1,D).
 \tag{P-F21b}
\]
Let $\Pi_D=\{\lambda_0v_0+\lambda_1v_1+\lambda_2v_2:
0\le\lambda_i<1\}$ be the half-open fundamental parallelepiped.
Every element of $S_D$ has a unique decomposition
\[
 v=p+k_0v_0+k_1v_1+k_2v_2,
 \qquad p\in\Pi_D\cap\mathbb Z^3,\quad k_i\in\mathbb Z_{\ge0}:
\]
write its unique nonnegative real ray coordinates and subtract their
integer parts. Conversely all such sums belong to $S_D$.
Therefore $R_D$ is a finite free module, with basis $X^p$, over the
polynomial subring
\[
 A_D=\mathbb C[X^{v_0},X^{v_1},X^{v_2}].
 \tag{P-F21c}
\]
The three displayed monomials are algebraically independent because
the rays are linearly independent. This proves finite generation,
Krull dimension three, and that they form a homogeneous system of
parameters of degrees $(1,1,D)$. It also proves Cohen--Macaulayness:
finite freeness over the polynomial subring gives flat localizations
with zero-dimensional fibers, so the localized rings are
Cohen--Macaulay; in particular the displayed parameter sequence
is regular.

We compute the free-module numerator explicitly. For a lattice point
$p\in\Pi_D$, its first two coordinates satisfy
\[
 x\in\{0,\ldots,D-1\},\quad
 y\in\{0,\ldots,D-2\},\quad
 \lambda_1=x/D,\quad\lambda_2=y/(D-1).
\]
There is a unique $\lambda_0\in[0,1)$ making the third coordinate
integral, and that coordinate is
\[
 s=y+\left\lceil\frac{x}{D}+\frac{y}{D-1}\right\rceil.
\]
For $y=0$, the point $x=0$ has degree zero and the other $D-1$
points have degree one. For $1\le y\le D-2$, the $D-y$ choices
$0\le x\le D-y-1$ have degree $y+1$, while the remaining $y$
choices have degree $y+2$. Indeed $x/D+y/(D-1)\le1$ is equivalent
to $x\le D-y-y/(D-1)$, whose integer upper bound is $D-y-1$.
It follows, including the empty-sum case $D=2$, that
\begin{align*}
 \sum_{p\in\Pi_D\cap\mathbb Z^3}t^{\deg p}
 &=1+(D-1)t+
       \sum_{y=1}^{D-2}\big((D-y)t^{y+1}+yt^{y+2}\big)\\
 &=1+(D-1)(t+\cdots+t^{D-1})+(D-2)t^D\\
 &=(1+(D-2)t)(1+t+\cdots+t^{D-1}).
\end{align*}
Dividing by $(1-t)^2(1-t^D)$, the Hilbert series of $A_D$, proves
the series in (P-F21a). Taking coefficients yields
\[
 f_D(s)=\binom{s+2}{2}+(D-2)\binom{s+1}{2},
\]
which is the stated polynomial for every $s\ge0$. All its ordinary
polynomial coefficients are positive. It is positive at every
nonnegative integer, so both stated indices are zero under the
definitions used above.

For completeness, normality can be checked without a presentation
of $R_D$. The cone is cut out by the integral linear inequalities
\[
 x\ge0,\qquad y\ge0,\qquad
 D(D-1)s-(D-1)x-D^2y\ge0.
\]
For each of these linear forms $\ell$, divide out the gcd of its
coefficients and form
$B_\ell=\mathbb C[\{v\in\mathbb Z^3:\ell(v)\ge0\}]$.
After an integral basis change, $B_\ell$ is a polynomial ring in
one variable over a Laurent polynomial ring in two variables, so
it is integrally closed. Inside the common Laurent polynomial
ring one has
$R_D=\bigcap_\ell B_\ell$, by comparing the supports of Laurent
polynomials. If an element of $\operatorname{Frac}(R_D)$ is integral
over $R_D$, it is integral over every $B_\ell$; since it lies in
the common rational function field, it belongs to each $B_\ell$
and hence to $R_D$. Thus $R_D$ is normal. It is an affine normal
toric ring, so it has rational singularities: Danilov,
\emph{The geometry of toric varieties}, Section~8.5,
Proposition~8.5.1, printed pp.~124--125, proves that a toric
subdivision $\rho:X'\to X$ has
$\rho_*\mathcal O_{X'}=\mathcal O_X$ and
$R^i\rho_*\mathcal O_{X'}=0$ for $i>0$; applying his toric
resolution in Section~8.2 gives the claim. The primary scan is
\url{https://www.math.purdue.edu/~arapura/danilov.pdf}.

Let $J_D$ be the span of all cone monomials not on the extreme ray
$\mathbb R_{\ge0}v_2$. This is an ideal: the nonnegative linear
form $\ell(x,y,s)=(D-1)s-Dy$ vanishes precisely on that ray,
so adding a cone element to an element off the ray stays off it.
Since $v_2$ is primitive, $S_D\cap\mathbb R_{\ge0}v_2
=\mathbb Z_{\ge0}v_2$. Thus
\[
 R_D/J_D\cong\mathbb C[z],\qquad z=X^{v_2},\quad\deg z=D.
 \tag{P-F21d}
\]
For any homogeneous parameter sequence $q_1,q_2,q_3$, the graded
ring $R_D/(q_1,q_2,q_3)$ is finite-dimensional. Some power of
$z$ therefore lies in $(q_1,q_2,q_3)$. In the quotient (P-F21d),
this is impossible if every $q_i$ maps to zero. Some parameter
has nonzero image in $\mathbb C[z]$, and its positive degree is
therefore a multiple of $D$. The tuple $(1,1,D)$ already
constructed is lexicographically minimal among nondecreasing
parameter-degree tuples: all positive degrees are at least one,
and every tuple with first two entries equal to one has last entry
at least $D$. This proves the exact canonical maximum.

Finally the quotient (P-F21d) has Hilbert series $1/(1-t^D)$.
The exact sequence $0\to J_D\to R_D\to\mathbb C[z]\to0$
shows that the least common multiple of the Hilbert denominators
of graded ideals of $R_D$ contains $\Phi_D$. By (P-F21a), the
ring's own reduced denominator is $(1-t)^3$: its numerator
evaluates to $D-1\ne0$ at one. This completes the proof.
\end{proof}

\paragraph{Prior example and precise consequence.}
The period-collapse triangle is classical. McAllister--Woods,
\emph{The minimum period of the Ehrhart quasi-polynomial of a
rational polytope}, arXiv \texttt{math/0310255}, Theorem~2.2 on
printed p.~3 and its proof on pp.~4--5, use
$\operatorname{conv}\{(0,0),(D,0),(1,(D-1)/D)\}$ and obtain
the polynomial in (P-F21a). The integral affine map
$(x,y)\mapsto(D-x-Dy,y)$ carries $P_D$ to that triangle;
on the $s$th dilation use $(x,y)\mapsto(Ds-x-Dy,y)$, which
is a lattice bijection. The source is
\url{https://arxiv.org/pdf/math/0310255}.
The algebraic consequences above are proved here without a claim
of historical priority.

Choosing $D=2^a$ makes the hidden universal factor a prime-power
factor of arbitrarily large order. The Krull dimension remains
three, the positive-form weights remain one, the numerator has
degree at most one and coefficient bit length $O(\log D)$, and
both indices remain zero. Hence normality, Cohen--Macaulayness,
rational singularities, positive Hilbert numerators, and global
quasi-polynomiality together do not give even a qualitative bound
on canonical parameter degrees in terms of the dimension and
positive-form weights. This rules out such a general transfer
as a way of strengthening P-F19. An argument for the weaker
ordinary plethysm PH3 must use additional structure of the actual
plethysm algebra.
